%============================================================
% SISTEMATIKA PENULISAN BUKU KP (5 Bab)
% Di-input oleh chapters/bab1.tex via \input{\SistematikaFile}
%============================================================

\vspace{0.3cm}
\begin{description}[leftmargin=0pt, labelsep=0pt, itemsep=4pt]
  \item[] \textbf{BAB 1}\quad \textbf{Pendahuluan}\\
    Bab ini menjelaskan latar belakang, rumusan masalah, tujuan, manfaat,
    ruang lingkup, metode kerja, dan susunan laporan.



  \item[] \textbf{BAB 2}\quad \textbf{Gambaran Perusahaan}\\
    Bab ini menjelaskan profil \NamaPerusahaan\ sebagai tempat pelaksanaan
    kerja praktik. Uraian mencakup struktur organisasi perusahaan, pembagian
    divisi kerja, hak dan wewenang tiap jabatan penting, lokasi perusahaan,
    penerapan K3 dan etika profesi, serta konteks penempatan pada Divisi Maintenance.

  \item[] \textbf{BAB 3}\quad \textbf{Landasan Teori}\\
    Bab ini berisi landasan teoritis yang menjadi acuan dalam pelaksanaan kerja
    praktik. Pembahasan mencakup sensor, RS485/Modbus RTU, Raspberry Pi dan
    HMI, SQLite, MQTT, dashboard server, besaran kelistrikan panel, pengolahan pada node, alarm, dan evaluasi fungsional.

  \item[] \textbf{BAB 4}\quad \textbf{Rancangan dan Implementasi Sistem}\\
    Bab ini menjelaskan proses perancangan dan implementasi sistem selama kerja
    praktik. Alur pembahasan mencakup observasi dan kebutuhan monitoring,
    pengembangan mesin cutting dan monitoring panel Power Room, integrasi
    dashboard server, pengujian dan evaluasi, dokumentasi hasil pekerjaan,
    kontribusi penulis, serta korelasi kegiatan dengan mata kuliah.

  \item[] \textbf{BAB 5}\quad \textbf{Penutup}\\
    Bab terakhir berisi kesimpulan dari kegiatan kerja praktik yang telah
    dilaksanakan serta hasil yang diperoleh. Selain itu, diberikan pula saran
    sebagai bahan masukan bagi perusahaan, institusi pendidikan, dan mahasiswa
    untuk pengembangan selanjutnya.
\end{description}
